\section{La firme intensive en connaissances et son modèle d'apprentissage}

Le chapitre précédent a décrit une situation de travail~; il reste à la doter des
instruments qui permettent de l'analyser. Trois corps de littérature sont mobilisés
successivement. Le premier caractérise ce qu'est structurellement un cabinet de conseil et
comment s'y forme un professionnel. Le deuxième explique sur quoi repose l'autorité de ce
professionnel devant son client. Le troisième examine ce que l'introduction d'une
technologie générative déplace dans cet édifice, en s'appuyant sur les travaux empiriques
les plus proches du terrain étudié. L'ordre n'est pas indifférent~: la question de
l'intelligence artificielle n'est traitée qu'en dernier, une fois établi ce qu'elle vient
percuter.

\subsection{Ce que produit une organisation dont le capital est le savoir}

Le concept de firme intensive en connaissances (\textit{knowledge-intensive firm}) désigne
les organisations dont l'activité repose sur un travail intellectuel complexe conduit par
une main-d'{\oe}uvre très qualifiée, et dont le capital réside dans l'expertise de ses
membres bien plus que dans ses actifs matériels \cite{starbuck1992}. Contrairement à
l'entreprise industrielle, où le capital est constitué de machines et d'installations, le
cabinet de conseil ne possède que ce que ses collaborateurs savent faire. Ce capital sort
de l'immeuble chaque soir, ce qui suffit à expliquer l'attention portée par ces firmes à la
rétention, à la formation et à la codification interne des savoirs.

La taxonomie des \textit{professional service firms} proposée par von Nordenflycht
\cite{nordenflycht2010} précise cette caractérisation par trois traits distinctifs~:
l'intensité de savoir, la faible intensité capitalistique et le degré de
professionnalisation de la main-d'{\oe}uvre. Le conseil en management occupe dans cette
taxonomie une position particulière. Il partage avec la médecine ou le droit l'intensité de
savoir et la faible intensité capitalistique, mais il ne dispose ni d'un ordre
professionnel, ni d'un diplôme d'exercice obligatoire, ni d'un code déontologique
opposable. Cette position intermédiaire a une conséquence directe sur le raisonnement
conduit dans ce mémoire~: faute de barrière institutionnelle à l'entrée, la légitimité du
consultant se rejoue à chaque mission plutôt qu'elle ne se déduit d'un titre. Elle autorise
par ailleurs, avec prudence, à transposer au conseil des résultats obtenus sur d'autres
professions intellectuelles, à condition de tenir compte de cet écart de
professionnalisation.

Alvesson \cite{alvesson2004} spécifie l'analyse au cas des cabinets et introduit la notion
décisive pour la suite~: l'ambiguïté de l'\textit{output}. Le livrable d'une mission de
conseil n'est pas mesurable comme l'est une pièce usinée. Sa qualité ne se constate pas de
façon indiscutable, ni au moment de la remise, ni souvent après. Le client ne peut
généralement pas vérifier par lui-même la justesse d'un diagnostic~; s'il le pouvait, il
n'aurait pas acheté la prestation. Il en résulte que le jugement porté sur le travail de
connaissance emprunte largement à des signaux indirects~: la réputation du cabinet, la
qualité formelle du document, l'assurance de celui qui le présente. Alvesson
\cite{alvesson2001} pousse ce constat plus loin en soutenant que le travail de connaissance
se juge sur l'image et l'identité autant que sur le résultat. Cette proposition, qui paraît
sévère, est en réalité l'articulation centrale du présent chapitre~: si la valeur perçue
tient pour partie à des signaux, alors toute technologie qui modifie la production de ces
signaux — la vitesse de rédaction, la densité apparente d'un document, la couverture
apparente d'un sujet — touche au c{\oe}ur de la prestation et non à sa périphérie.

\subsection{La bureaucratie professionnelle et la standardisation des qualifications}

La typologie des configurations organisationnelles proposée par Mintzberg
\cite{mintzberg1982} fournit la grille structurelle qui manque à ce qui précède. Le cabinet
de conseil s'apparente à une \textit{bureaucratie professionnelle}. Dans cette
configuration, le centre opérationnel est constitué de professionnels hautement qualifiés
qui disposent d'une large autonomie dans la conduite de leur travail. La technostructure y
est faible et la ligne hiérarchique mince, parce que l'essentiel de la coordination ne
passe ni par la supervision directe ni par la standardisation des procédés.
% Pagination à compléter par Arthur sur l'exemplaire consulté :
% \cite[p.~309]{mintzberg1982} pour le chapitre consacré à la bureaucratie professionnelle.

Le mécanisme de coordination propre à cette configuration est la \textit{standardisation des
qualifications}. L'organisation ne prescrit pas comment le travail doit être exécuté~; elle
s'assure en amont que celui qui l'exécute a été formé à le faire correctement. Le contrôle
est déplacé de l'acte vers la personne, et de l'exécution vers la formation. C'est là que le
modèle du conseil se distingue de la bureaucratie professionnelle canonique. Dans la
médecine ou l'enseignement, la standardisation des qualifications est assurée par des
institutions extérieures~: universités, ordres, examens d'État. Dans le conseil, ainsi que
le laissait prévoir la taxonomie de von Nordenflycht \cite{nordenflycht2010}, elle est
largement produite \textit{par le cabinet lui-même}. Le grade, la revue par le manager,
l'exposition graduée à la relation client et le passage par des missions de complexité
croissante remplacent le diplôme d'exercice.

Cette observation a une portée analytique forte. Elle signifie que le dispositif de
formation interne d'un cabinet n'est pas une fonction support parmi d'autres~: il est le
mécanisme de coordination principal de l'organisation. Une perturbation du parcours
d'apprentissage des juniors n'est donc pas un problème de gestion des ressources humaines
isolé, mais une atteinte à ce qui tient la structure ensemble. La question directrice
relative à l'impact de l'automatisation sur la montée en compétences des juniors devient,
sous cette grille, une question de théorie des organisations autant que de GRH.

\subsection{La conversion des connaissances et l'apprentissage par la pratique}

Reste à préciser \textit{comment} cette qualification se forme. La distinction établie par
Polanyi \cite{polanyi1966} entre savoir tacite et savoir explicite en fournit le point de
départ. Le savoir explicite est formulable, transmissible par le langage et stockable dans
un document~: une méthodologie, une norme, un modèle de livrable. Le savoir tacite est
incorporé dans la pratique et résiste à la mise en mots. Il regroupe le jugement, le sens
de la pertinence, la capacité à sentir qu'un chiffre est douteux ou qu'une recommandation
ne passera pas devant un comité. C'est précisément cette part du savoir professionnel qui
distingue un consultant expérimenté d'un débutant disposant des mêmes documents.

Nonaka \cite{nonaka1994}, puis Nonaka et Takeuchi \cite{nonaka1997}, opérationnalisent
cette distinction en un modèle de conversion des connaissances articulant quatre modes~:
socialisation (du tacite vers le tacite), externalisation (du tacite vers l'explicite),
combinaison (de l'explicite vers l'explicite) et internalisation (de l'explicite vers le
tacite). L'intérêt de ce modèle pour le présent travail tient à ce qu'il permet de nommer
précisément ce qui est en jeu. Deux modes seulement sont directement concernés par
l'automatisation de la production documentaire.

La \textit{socialisation} d'abord. Elle se joue dans le partage d'expérience directe~: le
junior qui observe un senior conduire un entretien client, qui reprend un livrable annoté,
qui entend un manager expliquer pourquoi telle formulation a été écartée. Elle ne passe pas
par le document mais par la coprésence dans l'activité. L'\textit{internalisation} ensuite~:
le savoir explicite ne devient jugement personnel qu'au prix d'une pratique répétée,
c'est-à-dire de l'apprentissage par la pratique (\textit{learning by doing}) au sens
ordinaire du terme.

C'est ici que se comprend la fonction réelle des tâches traditionnellement confiées aux
consultants débutants. Ces tâches sont bien identifiées~:
\begin{itemize}
    \item la recherche documentaire et la veille réglementaire~;
    \item la collecte, le tri et le traitement des données brutes de mission~;
    \item la première synthèse de rapports et la rédaction de livrables intermédiaires.
\end{itemize}

Perçues comme chronophages et à faible valeur ajoutée immédiate, elles remplissent en
réalité une fonction pédagogique critique. Elles obligent le débutant à se confronter à la
matière première désordonnée avant qu'elle ne soit mise en forme, et c'est cette
confrontation qui construit le jugement. Le junior qui a dû lire lui-même un texte
réglementaire pour en extraire l'exigence pertinente acquiert quelque chose que ne lui
donnera pas la lecture d'une synthèse déjà faite~: la connaissance de ce que le texte ne
dit pas, de ses ambiguïtés, des endroits où l'interprétation est ouverte. En termes du
modèle SECI, il ne se contente pas de recevoir de l'explicite~: il internalise.

L'hypothèse théorique que ce chapitre soumet à l'épreuve du terrain peut dès lors être
formulée sans ambiguïté. Un outil capable de produire instantanément la synthèse, le
premier jet du livrable ou le tableau comparatif ne supprime pas une tâche à faible valeur
ajoutée~: il supprime le support matériel d'un mode de conversion des connaissances. Le
livrable produit par la machine court-circuite l'internalisation, et le temps ainsi libéré
n'est réinvesti en socialisation que si l'organisation l'organise délibérément. Rien ne
garantit qu'elle le fasse.

\section{La fabrique de la légitimité de l'expert}

La section précédente a établi ce qu'un cabinet produit et comment il forme ceux qui
produisent. Elle laisse ouverte la question de savoir pourquoi un client accepte de payer,
puis de suivre, l'avis d'un professionnel dont il ne peut vérifier le travail. Cette
question relève de la sociologie des professions et de la théorie de la légitimité.

\subsection{La juridiction professionnelle et ses frontières}

Abbott \cite{abbott1988} propose de considérer les professions non comme des entités
isolées mais comme un système en concurrence permanente. Une profession n'existe que par la
\textit{juridiction} qu'elle parvient à revendiquer et à faire reconnaître, c'est-à-dire par
le lien exclusif qu'elle établit entre un domaine de problèmes et son propre savoir
abstrait. Ce lien s'établit par trois opérations caractéristiques~: le diagnostic, qui
classe un cas dans les catégories de la profession~; l'inférence, qui raisonne sur ce cas à
partir du savoir professionnel~; et le traitement, qui prescrit une action. La juridiction
n'est jamais acquise~: elle est disputée, devant les tribunaux, devant l'opinion et sur le
lieu de travail, par d'autres professions qui revendiquent le même territoire.

L'apport décisif de ce cadre est de désigner l'inférence comme le maillon vulnérable.
Abbott observe que les professions dont le diagnostic ou le traitement sont routinisables
perdent progressivement du terrain, tandis que la maîtrise de l'inférence — le raisonnement
sur les cas difficiles, ceux que la routine ne couvre pas — constitue le noyau défendable
d'une juridiction. Transposé au conseil, ce résultat oriente directement la troisième
question directrice. Si une technologie prend en charge une part croissante de la collecte
et de la mise en forme, la juridiction du consultant n'est pas nécessairement menacée~; elle
est déplacée vers ce qui reste non routinisable. Encore faut-il que le cabinet sache
identifier et revendiquer ce noyau, et que le client y reconnaisse une valeur. La question
n'est pas seulement de savoir ce que la machine sait faire, mais quel territoire la
profession parvient à redéfinir comme sien.

\subsection{Les trois registres de la légitimité}

Suchman \cite{suchman1995} fournit la grille permettant de décomposer ce que recouvre la
confiance accordée à un expert. La légitimité y est définie comme la perception
généralisée que les actions d'une entité sont désirables ou appropriées au regard d'un
système socialement construit de normes et de croyances. Trois registres sont distingués.
% Pagination à compléter par Arthur : tableau des registres dans \cite{suchman1995}.

La \textit{légitimité pragmatique} repose sur l'intérêt direct de l'audience~: le client
accorde du crédit parce qu'il en retire un bénéfice tangible, un problème résolu, un délai
tenu, une conformité obtenue. La \textit{légitimité cognitive} repose sur la
compréhensibilité et sur le caractère allant de soi de l'entité~: le consultant est crédible
parce que ce qu'il dit s'inscrit dans un cadre de référence partagé, parce que sa démarche
est intelligible et parce que le recours à un cabinet dans une telle situation va de soi. La
\textit{légitimité morale} repose sur un jugement normatif~: l'entité est légitime parce que
sa conduite est jugée conforme à ce qui doit être fait, indépendamment du bénéfice retiré.

Cette décomposition est opératoire pour le terrain étudié, car elle laisse prévoir que les
trois registres ne sont pas affectés de la même manière par l'automatisation. La légitimité
pragmatique peut se trouver renforcée~: un livrable plus rapide et plus complet sert
l'intérêt immédiat du client. La légitimité cognitive est ambivalente~: la production
accélérée de documents formellement irréprochables peut soutenir l'impression de maîtrise,
mais l'uniformisation des productions peut aussi éroder le caractère distinctif de
l'expertise. La légitimité morale est la plus exposée, parce qu'elle engage la question de
ce qui a été réellement fait et par qui. Un client qui découvre qu'une part substantielle de
l'analyse facturée a été produite sans intervention humaine significative ne conteste pas
d'abord la qualité du résultat~: il conteste la conformité de la conduite du cabinet à ce
qu'il croyait acheter.

\subsection{Ce sur quoi repose la confiance dans une prestation intellectuelle}

La jonction entre ces deux cadres et la section 2.1 s'opère par l'ambiguïté de l'output déjà relevée.
Si le résultat d'une mission de conseil n'est pas
directement évaluable, alors la légitimité n'est pas un supplément d'âme ajouté à la
prestation~: elle en constitue une part du produit lui-même \cite{alvesson2001}. Le client
n'achète pas seulement un diagnostic~; il achète l'assurance que ce diagnostic a été produit
par quelqu'un qui en répond.

Cette formulation permet de préciser la question centrale du mémoire. Ce qui est en jeu
lorsqu'une machine produit une part substantielle de la valeur intellectuelle n'est pas la
disparition de l'expert, mais le déplacement du fondement de son autorité. L'autorité
adossée au volume de travail fourni et à la maîtrise technique de la collecte s'affaiblit
mécaniquement. Celle qui s'adosse à la responsabilité assumée sur un jugement — au sens de
l'inférence d'Abbott et de la légitimité morale de Suchman — devient, en théorie, le
territoire défendable. Que ce déplacement soit effectivement perçu, revendiqué et valorisé
par les consultants et par leurs clients relève du terrain, et non de la théorie.

\section{L'intelligence artificielle générative comme technologie de rupture dans les
systèmes d'information}

Les deux sections précédentes ont construit la structure et la légitimité. Il reste à
caractériser la technologie qui les percute, sous un angle strictement organisationnel~: ce
qui importe ici n'est pas la manière dont ces outils fonctionnent, mais la manière dont ils
sont adoptés, ce qu'ils déplacent dans la division du travail et ce qu'ils font au jugement
professionnel.

\subsection{Les modèles d'acceptation technologique et leurs limites}

Le modèle d'acceptation de la technologie (\textit{Technology Acceptance Model}) formulé par
Davis \cite{davis1989} explique l'usage effectif d'un système d'information par deux
variables perçues~: l'utilité perçue et la facilité d'utilisation perçue. Sa longévité tient
à sa parcimonie. Ses prolongements successifs en ont progressivement corrigé le caractère
individualiste. Le TAM2 \cite{venkatesh2000} introduit l'influence sociale et les processus
instrumentaux cognitifs~: l'usage d'un pair, la norme subjective et l'image tirée de
l'adoption pèsent sur la décision d'utiliser. Le modèle unifié UTAUT \cite{venkatesh2003}
ajoute aux déterminants de l'intention les \textit{conditions facilitatrices},
c'est-à-dire l'ensemble des ressources organisationnelles rendant l'usage possible. Le TAM3
\cite{venkatesh2008} franchit un pas supplémentaire en identifiant les interventions
organisationnelles susceptibles d'agir sur l'acceptation~: soutien managérial, formation,
participation des utilisateurs.

Ces prolongements servent la première et la quatrième question directrices. Le TAM2 éclaire
la diffusion par imitation entre pairs, plausible dans des équipes projet où l'usage d'un
outil se transmet par observation plutôt que par prescription. L'UTAUT désigne le
management comme producteur des conditions facilitatrices, ce qui est exactement l'objet de
la question de gouvernance. Le TAM3 fournit le répertoire des leviers actionnables et
constitue à ce titre l'appui théorique des recommandations formulées ultérieurement.

Deux réserves doivent cependant être posées. La première tient à l'objet~: ces modèles ont
été construits sur des logiciels de bureautique dont l'usage était prescrit par
l'employeur, non sur des outils que le collaborateur peut mobiliser de sa propre
initiative, hors du cadre officiel, y compris depuis un terminal personnel. La distinction
entre adoption prescrite et appropriation spontanée, que le TAM ne thématise pas, est
centrale sur le terrain étudié. La seconde tient à la variable expliquée~: ces modèles
expliquent l'intention d'usage, non les effets de l'usage sur la compétence de
l'utilisateur. Or c'est précisément l'effet, et non l'adoption, qui est ici en question. Le
cadre de l'acceptation est donc nécessaire mais insuffisant~; il doit être complété par les
travaux empiriques portant sur les conséquences de l'usage.

\subsection{Du consultant augmenté au consultant assisté~: ce qu'établissent les travaux
empiriques}

L'expérimentation la plus proche du terrain de ce mémoire a été conduite en 2023 auprès de
consultants du Boston Consulting Group, en collaboration avec le centre de recherche du
groupe, et publiée depuis dans \textit{Organization Science}
\cite{dellacqua2026}. Son résultat central est la notion de \textit{frontière dentelée}
(\textit{jagged frontier})~: le périmètre des tâches sur lesquelles l'outil apporte un gain
n'est ni continu ni prévisible depuis l'extérieur. Sur les tâches situées à l'intérieur de
ce périmètre, les consultants équipés produisent davantage, plus vite et avec une qualité
supérieure. Sur des tâches d'apparence voisine mais situées au-delà de la frontière, l'usage
de l'outil dégrade la performance~: les consultants assistés se trompent plus souvent que
ceux qui travaillent sans assistance.

Ce résultat est le plus important du chapitre, pour deux raisons. D'abord parce qu'il
interdit de poser la question en termes de gain global~: l'effet est conditionnel à la
nature de la tâche, et la difficulté pratique tient à ce que la frontière est invisible pour
celui qui l'approche. Ensuite parce qu'il déplace le problème managérial. Si la frontière
était identifiable, il suffirait de prescrire l'usage en deçà et de l'interdire au-delà.
Comme elle ne l'est pas, la seule ressource disponible est le jugement du professionnel sur
sa propre situation — c'est-à-dire exactement la compétence dont la section 2.1 laisse
craindre l'érosion.

Deux travaux complémentaires portent sur la distribution des gains. Noy et Zhang
\cite{noy2023} établissent, sur des tâches de rédaction professionnelle, un gain de
productivité assorti d'un resserrement de l'écart entre les participants les moins
performants et les plus performants~: l'outil profite davantage aux premiers. Brynjolfsson,
Li et Raymond \cite{brynjolfsson2025} retrouvent ce schéma en conditions réelles de
déploiement en entreprise, où les gains se concentrent sur les travailleurs les moins
expérimentés. La convergence entre un protocole expérimental et des données d'entreprise
donne à ce résultat une robustesse peu commune dans la littérature récente.

Sa lecture, en revanche, est ambivalente, et cette ambivalence structure la deuxième
question directrice. Une lecture optimiste y voit une accélération de la montée en
compétences~: le junior atteint plus tôt un niveau de production acceptable. Une lecture
pessimiste observe que le resserrement porte sur la production, non sur la compétence~: si
le débutant atteint le résultat du senior sans avoir parcouru le chemin qui produit le
jugement, l'écart de performance mesurée masque un écart de compétence inchangé, voire
creusé. La théorie ne permet pas de trancher entre ces deux lectures~; le terrain, oui.
C'est précisément ce que le dispositif d'entretiens du chapitre suivant cherche à établir en
confrontant le ressenti des juniors sur leur propre apprentissage à la perception qu'en ont
leurs encadrants.

Enfin, les travaux d'Eloundou, Manning, Mishkin et Rock \cite{eloundou2024} sur l'exposition
des métiers aux modèles de langage situent ce mouvement à l'échelle macro-économique et
confirment que les professions intellectuelles à forte composante rédactionnelle et
analytique figurent parmi les plus exposées. Cette référence sert au cadrage général~; elle
ne dit rien de ce qui se joue à l'échelle d'une équipe.

\subsection{Opacité, contrôle et atrophie de la réflexivité critique}

Le dernier ensemble de travaux porte sur ce que l'usage d'une technologie opaque fait au
jugement professionnel. Anthony \cite{anthony2021} conduit une étude ethnographique
d'analystes juniors confrontés à une technologie analytique dont le fonctionnement leur
échappe. Le résultat est nuancé et d'autant plus utile~: le clivage ne passe pas entre ceux
qui utilisent l'outil et ceux qui s'en abstiennent, mais entre deux manières de l'utiliser.
Les analystes qui interrogent la production de l'outil, qui cherchent à comprendre d'où
vient un résultat et à en éprouver la cohérence, développent une expertise plus solide. Ceux
qui l'utilisent sans l'interroger — parce que le résultat paraît plausible, parce que le
délai presse, parce que l'outil fait autorité — développent une expertise plus
superficielle. Le mécanisme est identifié~: c'est l'acceptation non critique du résultat, et
non l'usage lui-même, qui appauvrit l'apprentissage.

Ce résultat est le plus directement transposable au cas étudié. Il désigne une variable
managériale actionnable, là où le constat d'une atrophie générale des compétences serait
resté sans prise. Il fournit également un critère d'analyse pour le chapitre~4~: la question
posée aux répondants n'est pas de savoir s'ils utilisent ces outils, mais ce qu'ils font du
résultat produit, et à quelles conditions ils le remettent en cause.

Lebovitz, Levina et Lifshitz-Assaf \cite{lebovitz2021} apportent un argument de fond
souvent absent des discussions sur la fiabilité de ces outils. Les « vérités terrain » sur
lesquelles une technologie d'intelligence artificielle est entraînée et évaluée sont
elles-mêmes construites et contestables~; l'évaluation d'un outil suppose donc un étalon qui
n'est pas toujours disponible. Appliquée au conseil, la portée de cet argument est
considérable~: le problème posé par les productions erronées n'est pas seulement que l'outil
se trompe, c'est que le professionnel ne dispose pas toujours d'une référence indiscutable
contre laquelle vérifier. Le contrôle qualité ne peut donc être conçu comme une simple
opération de vérification factuelle.

Les mêmes auteurs \cite{lebovitz2022} examinent comment des professionnels arbitrent
concrètement face à cette opacité lorsqu'ils doivent porter un jugement engageant leur
responsabilité. Le parallèle avec le consultant qui signe un livrable remis au client est
direct. Cette étude constitue l'appui principal de la quatrième question directrice et la
matrice des protocoles de validation qui seront proposés en recommandation.

Kellogg, Valentine et Christin \cite{kellogg2020} complètent le dispositif en analysant le
contrôle algorithmique comme un nouveau terrain de lutte dans l'organisation. Les auteurs
recensent les leviers par lesquels une technologie oriente, évalue et discipline le travail,
ainsi que les formes de résistance que ces leviers suscitent. Cette grille est
indispensable à la formulation des recommandations de gouvernance, car elle rappelle
qu'encadrer un usage n'est jamais un acte neutre~: toute règle produit des effets
d'appropriation et de contournement. Elle éclaire directement la tension centrale de la
quatrième question directrice, entre la nécessité de maîtriser un risque opérationnel et le
maintien de l'autonomie professionnelle sans laquelle une bureaucratie professionnelle ne
fonctionne pas.

\subsection{Synthèse~: quatre questionnements, quatre cadres d'analyse}

L'ensemble des cadres mobilisés peut être rapporté aux quatre questionnements directeurs du
mémoire, chacun trouvant ici les instruments de son traitement empirique.

\begin{table}[h]
\centering
\small
\begin{tabular}{|p{0.24\linewidth}|p{0.33\linewidth}|p{0.33\linewidth}|}
\hline
\textbf{Questionnement} & \textbf{Cadre théorique mobilisé} & \textbf{Appui empirique} \\
\hline
1. Appropriation réelle des outils
& TAM et prolongements~: influence sociale, conditions facilitatrices
(\textsc{Davis}, 1989~; \textsc{Venkatesh} et \textsc{Davis}, 2000~;
\textsc{Venkatesh} \textit{et al.}, 2003)
& Distribution non uniforme des usages et des gains
(\textsc{Dell'Acqua} \textit{et al.}, 2026) \\
\hline
2. Apprentissage et montée en compétences des juniors
& Bureaucratie professionnelle et standardisation des qualifications
(\textsc{Mintzberg}, 1982)~; conversion des connaissances et savoir tacite
(\textsc{Polanyi}, 1966~; \textsc{Nonaka} et \textsc{Takeuchi}, 1997)
& Resserrement des écarts de performance (\textsc{Noy} et \textsc{Zhang}, 2023~;
\textsc{Brynjolfsson} \textit{et al.}, 2025)~; expertise superficielle en cas d'usage non
critique (\textsc{Anthony}, 2021) \\
\hline
3. Légitimité de l'expert face au client
& Juridiction professionnelle et inférence (\textsc{Abbott}, 1988)~; registres
pragmatique, cognitif et moral (\textsc{Suchman}, 1995)~; ambiguïté de l'\textit{output}
(\textsc{Alvesson}, 2001) & Déplacement de la valeur vers les tâches non routinisables
(\textsc{Dell'Acqua} \textit{et al.}, 2026) \\
\hline
4. Gouvernance et maîtrise des risques
& Contrôle algorithmique et résistances (\textsc{Kellogg} \textit{et al.}, 2020)~;
interventions organisationnelles sur l'acceptation (\textsc{Venkatesh} et \textsc{Bala},
2008) & Jugement professionnel en situation d'opacité (\textsc{Lebovitz} \textit{et al.},
2021 et 2022) \\
\hline
\end{tabular}
\caption{Correspondance entre les questionnements directeurs et les cadres d'analyse
mobilisés}
\label{tab:cadres-questionnements}
\end{table}

Le tableau~\ref{tab:cadres-questionnements} rapporte l'ensemble des cadres mobilisés aux quatre questionnements directeurs. Les cadres retenus ne sont pas juxtaposés~: ils se
répondent. La bureaucratie professionnelle explique pourquoi une perturbation de
l'apprentissage est une question structurelle et non une question de formation~; la théorie
de la légitimité explique pourquoi cette perturbation finit par atteindre la relation
client~; les travaux empiriques sur l'usage établissent que l'effet dépend de la manière
d'utiliser plus que de l'usage lui-même, ce qui redonne au management une prise. La chaîne
de raisonnement soumise au terrain peut être énoncée en une phrase~: l'automatisation des
tâches d'analyse déplace le lieu où se forme le jugement professionnel, et la question
managériale est de savoir si l'organisation reconstitue ailleurs ce qu'elle supprime là.

Le chapitre suivant expose la démarche retenue pour soumettre cette proposition à
l'épreuve des faits.